% !TeX root = ../概率论与数理统计笔记.tex
% 由主文档通过 \input 载入；请编译主文档。
\setcounter{chapter}{5} % 保留原书章号。
\chapter{随机变量函数的分布}
\label{chap:ps-transformations}

本章研究由已知随机量构造的新随机量，其分布如何求得。
先修内容为联合分布与微积分；学习目标是根据变换性质选择分布函数法、变量替换法、卷积或条件化。

\section{分布函数法与离散变换}
\label{sec:ps-cdf-transform}
从事件 $\{g(X)\leq y\}$ 出发求变换后的分布，明确函数 $g$、原变量 $X$ 与新变量 $Y=g(X)$ 的角色。
离散情形对相同函数值合并概率；连续情形处理单调变换、平方与绝对值等非单调变换。
优先确定新变量的取值范围，并保留可能存在的原子概率。

\section{密度变换与雅可比行列式}
\label{sec:ps-jacobian}
在适当的可微、可逆及雅可比非退化条件下建立一维和多维密度变换公式。
说明使用逆变换的\emph{雅可比行列式}（Jacobian determinant）绝对值，并将原支撑区域变换到新坐标中。
多分支情形分区求和；用极坐标和商的分布展示方法的适用范围。

\section{和、差及其他组合的分布}
\label{sec:ps-convolution}
建立独立随机变量和的\emph{卷积}（convolution）公式，比较离散求和与连续积分。
推导独立同成功概率二项变量、独立泊松变量、独立正态变量以及独立同率参数伽马变量的可加性。
对于相关变量，回到联合分布或条件分布计算，不能直接套用独立卷积。

\section{最大值、最小值与次序统计量的准备}
\label{sec:ps-extreme-transform}
对独立样本用事件交并求最大值和最小值的分布，比较同分布与不同分布情形。
介绍样本排序的含义，为第\ref{chap:ps-sampling}章的次序统计量作准备。
用极值与均值的不同归一化说明不同统计量可能具有不同的误差尺度。

\section{例题、习题与学习检查}
规划正态变量平方、独立均匀变量之和、两个独立指数变量的最小值、连续变量之比等例题。
每题按“支撑区域---变换关系---分布计算---归一化核对”组织。
习题安排同一问题的两种解法，并检查是否遗漏非单调变换的分支。






